\section{An exponential information gap for wide symplectic blocks}
\label{sec:symplectic-information}

The second result concerns a finite linear-algebra mechanism developed in the
source project.  A rectangular block of a uniformly random symplectic
transformation can have almost the same rank distribution as a matrix of
independent uniform entries, while the two complete matrix distributions are
almost mutually singular.  We determine both discrepancies and the number of
rank observations needed to distinguish the models.

The exact arbitrary-frame compression law and a square-block version of rank
blindness already appear in the supplied article
\cite[sections on symplectic compressions and Lagrangian projections]{ProveItArticle}.
The square result gives a positive limiting discrepancy between complete
matrices.  Here the number of columns exceeds the number of rows by a growing
amount.  In this regime the complete-matrix discrepancy tends to one, and we
prove a sharp rank-only sample scale.  The finite orbit count below is a
specialization of the article's law, for which we give a direct proof.  The
uniform asymptotic separation and statistical sample scale are the additional
conclusions.  They are unconditional finite-field theorems; they do not use
Dirichlet GRH.  Classical uniform-matrix rank distributions provide the
comparison law \cite{FulmanGoldstein2015}; prescribed-form sampling is also
studied in \cite{Ishizuka2026}.

\subsection{The two random-matrix models}

Fix a prime power $q$, integers
\begin{equation}
  n\ge c\ge1,\qquad b=n+c\text{ even},\qquad m=b/2,
  \qquad Q=q^{-1},\qquad h=n-c.
  \label{symp:eq:parameters}
\end{equation}
In particular $h$ is even and $0\le h\le n-1$.  Let $H$ be a
nondegenerate alternating $b\times b$ matrix over $\mathbb F_q$.  Here
``alternating'' means $x^{\mathsf T}Hx=0$ for every $x$; in characteristic two
this includes the zero-diagonal condition.  Set
\[
 J_n=\begin{pmatrix}0&I_n\\-I_n&0\end{pmatrix}.
\]
Choose a $2n\times b$ matrix $X_0$ satisfying
$X_0^{\mathsf T}J_nX_0=H$.  Such a matrix exists since $b\le2n$; its columns
are independent because $H$ is nondegenerate.  For uniform
$g\in\operatorname{Sp}_{2n}(\mathbb F_q)$, write
\begin{equation}
 gX_0=\begin{pmatrix}A\\B\end{pmatrix},
 \qquad A,B\in\operatorname{Mat}_{n,b}(\mathbb F_q).
 \label{symp:eq:model}
\end{equation}
Let $U$ be uniform on $\operatorname{Mat}_{n,b}(\mathbb F_q)$, so its entries
are independent and uniform.  The two laws are independent of the particular
choice of $X_0$ with the prescribed Gram matrix.

For probability measures on a finite set we use
\[
 \operatorname{TV}(P,V)=\frac12\sum_x|P(x)-V(x)|,
 \qquad
 \chi^2(P\Vert V)=\sum_x\frac{(P(x)-V(x))^2}{V(x)}
\]
when $V$ is positive on the support of $P$.  Write
$\operatorname{TV}(A,U)$ for the distance between their laws.  Define
\begin{equation}
 a_j=\prod_{i=1}^j(1-Q^i),\quad a_0=1,\quad
 a_\infty=\prod_{i=1}^\infty(1-Q^i)>0,\quad
 C=\prod_{j=n-m+1}^{n}(1-Q^{2j}).
 \label{symp:eq:products}
\end{equation}
All finite products with empty index sets are one.  For a proposed matrix
$M$, let $K=\ker M$ and $k=\dim K$.  Necessarily $c\le k\le b$.
We call $K$ \emph{$H$-isotropic} if $x^{\mathsf T}Hy=0$ for every $x,y\in K$.

\begin{theorem}[Exact law and optimal distinguishing events]
\label{symp:thm:finite}
In the model \eqref{symp:eq:model}, the likelihood relative to a uniform
matrix is
\begin{equation}
 \frac{\Pr(A=M)}{q^{-nb}}
 =\mathbf 1_{\{\ker M\text{ is }H\text{-isotropic}\}}
   \frac{q^{\binom{k}{2}}}{C}.
 \label{symp:eq:density}
\end{equation}
Consequently,
\begin{equation}
 \operatorname{TV}(A,U)
 =\Pr\bigl(H|_{\ker U\times\ker U}\ne0\bigr).
 \label{symp:eq:fulltv}
\end{equation}
For $c\le k\le b$ put
\begin{equation}
 u(k)=Q^{k(k-c)}\frac{a_na_b}{a_ka_{k-c}a_{b-k}},\qquad
 R_{b,k}=\prod_{i=0}^{k-1}
          \frac{1-Q^{b-2i}}{1-Q^{b-i}}\quad(k\le m),
 \label{symp:eq:rank-formulas}
\end{equation}
and set $R_{b,k}=0$ for $k>m$.  Then
\begin{equation}
 \Pr(\dim\ker U=k)=u(k),\qquad
 \Pr(\dim\ker A=k)=u(k)\frac{R_{b,k}}C.
 \label{symp:eq:rank-law}
\end{equation}
Full row rank is the only rank class with more probability under $A$ than
under $U$.  In particular, with
\[
 D=\prod_{j=m+1}^n(1-Q^{2j}),\qquad
 E=\prod_{j=n+1}^{n+c}(1-Q^j),
\]
the exact rank discrepancy is
\begin{equation}
 \Delta:=\operatorname{TV}(\operatorname{rank}A,\operatorname{rank}U)
 =\frac{a_n}{a_c}\bigl(D^{-1}-E\bigr).
 \label{symp:eq:rank-tv}
\end{equation}
The overlap of the complete matrix laws is
\begin{equation}
 \Omega:=1-\operatorname{TV}(A,U)
  =\sum_{k=c}^{m}u(k)Q^{\binom{k}{2}}R_{b,k}.
 \label{symp:eq:overlap}
\end{equation}
\end{theorem}

\subsection{Counting frames and fibres}

We first prove the likelihood without using the general compression formula.
The symplectic group acts transitively on the frames $X$ satisfying
$X^{\mathsf T}J_nX=H$: identify the nondegenerate spans of two frames by the
map taking corresponding columns to each other, and extend that isometry
over their symplectic orthogonal complements.  A fixed frame has stabilizer
$\operatorname{Sp}_{2(n-m)}(\mathbb F_q)$, acting on its complement.  The
standard count
\[
 |\operatorname{Sp}_{2t}(\mathbb F_q)|
   =q^{t^2}\prod_{j=1}^{t}(q^{2j}-1)
\]
therefore gives
\begin{equation}
 \#\{X:X^{\mathsf T}J_nX=H\}
 =q^{2nb-b(b-1)/2}C.
 \label{symp:eq:frame-count}
\end{equation}
The group-order formula itself follows by successively choosing a nonzero
first vector, a partner pairing to one with it, and a symplectic basis in
their orthogonal complement.

Fix $A=M$ of rank $r=b-k$.  The equation for the remaining rows is
\begin{equation}
 M^{\mathsf T}B-B^{\mathsf T}M=H.
 \label{symp:eq:gram-fibre}
\end{equation}
Its restriction to $K\times K$ is zero, so $H$-isotropy of $K$ is necessary.
Choose invertible matrices $P,T$ with
$PMT=\left(\begin{smallmatrix}I_r&0\\0&0\end{smallmatrix}\right)$.
Replacing $B$ by $P^{-\mathsf T}BT$ and $H$ by $T^{\mathsf T}HT$ preserves
\eqref{symp:eq:gram-fibre}.  Partition rows into sizes $r,n-r$ and columns
into sizes $r,k$.  In these coordinates the equation prescribes
\[
 B_{11}-B_{11}^{\mathsf T}=H_{11},\qquad B_{12}=H_{12},
 \qquad H_{22}=0.
\]
If $K$ is isotropic, the last equation holds.  The first has
$q^{r(r+1)/2}$ solutions, since the diagonal and one entry in each
off-diagonal pair are free.  The last $n-r$ rows of $B$ are entirely free.
Thus the number of solutions is
\begin{equation}
 q^{r(r+1)/2+(n-r)b}
  =q^{nb-b(b-1)/2+\binom{k}{2}}.
 \label{symp:eq:fibre-count}
\end{equation}
Every such solution gives an independent frame: a vector in the kernel of
$\left(\begin{smallmatrix}M\\B\end{smallmatrix}\right)$ belongs to
$\operatorname{rad}H=0$.  This argument requires no division by two and
applies in characteristic two.  Dividing \eqref{symp:eq:fibre-count} by
\eqref{symp:eq:frame-count} proves \eqref{symp:eq:density}.

Since $C<1$, every positive likelihood in \eqref{symp:eq:density} is
strictly larger than one.  The entire deficit relative to the uniform law
therefore lies outside the support, proving \eqref{symp:eq:fulltv}.  If $Z$
is a basis matrix for $\ker M$, the optimal event is simply
$Z^{\mathsf T}HZ\ne0$.  It can be evaluated using a polynomial number of
finite-field operations.

\subsection{The rank law and its sign pattern}

To count uniform matrices with kernel dimension $k$, choose a $k$-subspace
of $\mathbb F_q^b$ and an injective map from its quotient to
$\mathbb F_q^n$.  Dividing the resulting count by $q^{nb}$ gives $u(k)$
in \eqref{symp:eq:rank-formulas}.  Conditional on $k$, the kernel is a
uniform $k$-subspace.

For an independent isotropic $k$-frame in the nondegenerate $b$-space,
after $i$ vectors have been chosen the next vector must lie in their common
orthogonal complement and outside their span.  There are
$q^{b-i}-q^i$ choices.  Comparing with all independent $k$-frames gives
\begin{equation}
 \Pr(K\text{ is }H\text{-isotropic}\mid\dim K=k)
 =\frac{\prod_{i=0}^{k-1}(q^{b-i}-q^i)}
        {\prod_{i=0}^{k-1}(q^b-q^i)}
 =Q^{\binom{k}{2}}R_{b,k}
 \label{symp:eq:isotropy}
\end{equation}
for $k\le m$, and the probability is zero for $k>m$.  Averaging
\eqref{symp:eq:density} proves \eqref{symp:eq:rank-law}; summing
\eqref{symp:eq:isotropy} proves \eqref{symp:eq:overlap}.

It remains to identify the sign of every rank discrepancy.  Cancellation
of factors gives
\begin{equation}
 \frac{R_{b,c}}C=(DE)^{-1}>1.
 \label{symp:eq:fullrank-ratio}
\end{equation}
For $1\le k<m$ the successive ratio
\[
 \frac{R_{b,k+1}}{R_{b,k}}
  =\frac{1-Q^{b-2k}}{1-Q^{b-k}}
\]
is strictly smaller than one.  We prove that $R_{b,c+1}<C$, after which
all remaining inequalities follow.  If $c=n$, then $c=m$ and this is
immediate.  Otherwise $h=n-c\ge2$, and direct cancellation gives
\begin{equation}
 \frac{R_{b,c+1}}C
 =\frac{1-Q^{n-c}}
 {\displaystyle\prod_{j=n}^{n+c}(1-Q^j)
  \displaystyle\prod_{j=m+1}^{n}(1-Q^{2j})}.
 \label{symp:eq:first-deficiency}
\end{equation}
The sum of the quantities subtracted from one in this denominator, divided
by $Q^{n-c}$, is at most
\begin{equation}
 \frac{Q^c(1-Q^{c+1})}{1-Q}
 +\frac{Q^{2c+2}}{1-Q^2}<1.
 \label{symp:eq:sign-bound}
\end{equation}
Indeed, for $c=1$ the upper bound is at most
$3/4+1/12=5/6$; for $c\ge2$ it is at most
$1/2+1/48=25/48$.  We used only $Q\le1/2$.  The inequality
$\prod_i(1-x_i)\ge1-\sum_i x_i$ now shows that the denominator in
\eqref{symp:eq:first-deficiency} exceeds its numerator.  Hence full row
rank is the sole rank class with positive discrepancy.  Multiplying its
likelihood excess by $u(c)=a_b/a_c$ proves
\eqref{symp:eq:rank-tv}, completing the proof of
Theorem~\ref{symp:thm:finite}.

\subsection{Uniform separation as the excess width grows}

\begin{theorem}[Almost singular matrices with almost identical ranks]
\label{symp:thm:separation}
Fix $q$.  As $c\to\infty$, uniformly for every $n\ge c$ with $n+c$ even,
\begin{equation}
 \Delta=\frac{q^{-n}}{q-1}\bigl(1+O_q(q^{-c})\bigr).
 \label{symp:eq:rank-asymptotic}
\end{equation}
An explicit estimate, valid throughout \eqref{symp:eq:parameters}, is
\begin{equation}
 \left|(q-1)q^n\Delta-1\right|\le3q^{-c}.
 \label{symp:eq:rank-explicit}
\end{equation}
Define
\begin{align}
 \Omega_0&=q^{-\binom c2}\frac{a_n}{a_c}
  \prod_{j=(n-c)/2+1}^{(n+c)/2}(1-q^{-2j}),
  \label{symp:eq:leading-overlap}\\
 \eta_q(c)&=\frac{a_\infty^{-2}q^{-(2c+1)}}{1-q^{-(2c+4)}}.
  \label{symp:eq:overlap-error}
\end{align}
Then
\begin{equation}
 \Omega_0\le\Omega\le\Omega_0\bigl(1+\eta_q(c)\bigr).
 \label{symp:eq:overlap-certified}
\end{equation}
In particular, for the even integer $h=n-c$ put
\[
 B_{q,h}=\prod_{j=h/2+1}^{\infty}(1-q^{-2j}).
\]
Uniformly for $h\ge0$,
\begin{equation}
 1-\operatorname{TV}(A,U)
 =B_{q,h}\,q^{-\binom c2}\bigl(1+O_q(q^{-c})\bigr).
 \label{symp:eq:full-asymptotic}
\end{equation}
Thus $\operatorname{TV}(\operatorname{rank}A,\operatorname{rank}U)\to0$
and $\operatorname{TV}(A,U)\to1$.  If also $n-c\to\infty$, the leading
factor $B_{q,h}$ tends to one.  No such additional hypothesis is needed for
the separation itself.
\end{theorem}

\begin{proof}
For the rank estimate, write
\[
 S_E=\sum_{j=n+1}^{n+c}Q^j
     =\frac{Q^n(1-Q^c)}{q-1},\qquad
 S_D=\sum_{j=m+1}^{n}Q^{2j}
     \le\frac{Q^{n+c}}{q^2-1}.
\]
The elementary product estimates give
\begin{gather}
 S_E-\tfrac12S_E^2\le1-E\le S_E,\qquad
 0\le D^{-1}-1\le\frac{S_D}{1-S_D},
 \label{symp:eq:product-bounds}\\
 1-\frac{Q^c}{q-1}\le\frac{a_n}{a_c}\le1.
 \nonumber
\end{gather}
For the first lower bound, the sum of pairwise products is at most
$S_E^2/2$, and the two-term Bonferroni inequality applies.  Since
$n+c\ge2$, we have $S_D\le1/12$.  Divide
\eqref{symp:eq:rank-tv} by $Q^n/(q-1)$.  The upper bound from
\eqref{symp:eq:product-bounds} is at most
$1+(4/11)Q^c$.  The lower bound is at least
\[
 \left(1-\frac{Q^c}{q-1}\right)
 \left(1-Q^c-\frac{Q^n}{2(q-1)}\right)
 \ge1-\frac52Q^c.
\]
The weaker symmetric bound \eqref{symp:eq:rank-explicit} follows.

The $k=c$ summand of \eqref{symp:eq:overlap} is exactly
$\Omega_0$.  To bound the others, put $k=c+j$.  Since $R_{b,c+j}\le R_{b,c}$,
and
\begin{equation}
 \frac{u(c+j)}{u(c)}
  =Q^{j(c+j)}\frac{a_ca_n}{a_{c+j}a_ja_{n-j}}
  \le a_\infty^{-2}Q^{j(c+j)},
 \label{symp:eq:overlap-ratio}
\end{equation}
the ratio of the $j$th summand to $\Omega_0$ is at most
\begin{equation}
 a_\infty^{-2}Q^{2cj+(3j^2-j)/2}.
 \label{symp:eq:overlap-tail}
\end{equation}
For $j\ge1$, successive exponents increase by
$2c+3j+1\ge2c+4$.  A geometric majorant proves
\eqref{symp:eq:overlap-certified}.  Finally,
$a_n/a_c=1+O_q(Q^c)$, while the finite even product in
\eqref{symp:eq:leading-overlap} differs from $B_{q,h}$ by relative
$O_q(Q^{n+c})$.  Because $n+c\ge2c$, these estimates prove
\eqref{symp:eq:full-asymptotic}, uniformly in $h$.
\end{proof}

The boundary $c=n$ is included.  Here the source frame fills the ambient
space, so $A$ has full row rank surely.  The exact formulas reduce to
\begin{equation}
 \Delta=1-\prod_{j=n+1}^{2n}(1-q^{-j}),\qquad
 \Omega=q^{-\binom n2}\prod_{j=1}^{n}(1-q^{-2j}).
 \label{symp:eq:saturation}
\end{equation}
The second equality also counts the uniformly random matrices whose
$n$-dimensional kernels are Lagrangian.  More generally, when $h$ remains
fixed, the factor $B_{q,h}$ records the finite distance to this boundary.

\subsection{Sharp sample complexity when only ranks are observed}

Small one-observation total variation by itself does not identify the right
sample scale.  The following calculation does so.

\begin{theorem}[Rank-only information and its sharp sample scale]
\label{symp:thm:sample}
Let $P$ be the law of $\operatorname{rank}A$ and let $V$ be the law of
$\operatorname{rank}U$.  For fixed $q$ and $c\to\infty$, uniformly over
\eqref{symp:eq:parameters},
\begin{equation}
 \chi^2(P\Vert V)
  =\frac{q^{-(2n-c)}}{q-1}\bigl(1+O_q(q^{-c})\bigr).
 \label{symp:eq:chi-square}
\end{equation}
For $t$ independent matrix observations, with only their ranks retained,
the following conclusions hold.
\begin{enumerate}
\item If $t=o(q^{2n-c})$, the total variation distance between the two
 joint rank laws tends to zero.  Under equal prior probabilities, every
 rank-only decision rule consequently has error probability tending to
 $1/2$ or larger.
\item If $t/q^{2n-c}\to\infty$, a threshold on the number of matrices
 failing to have full row rank distinguishes the models with error
 probability tending to zero.
\end{enumerate}
For any fixed target error strictly between zero and $1/2$, the required
number of rank observations is therefore $\Theta_q(q^{2n-c})$, with the
implicit constants also allowed to depend on the target error.
By contrast, one complete matrix has optimal equal-prior error
\begin{equation}
 \frac{\Omega}{2}
 =\frac{B_{q,n-c}}2\,q^{-\binom c2}
      \bigl(1+O_q(q^{-c})\bigr).
 \label{symp:eq:one-matrix-risk}
\end{equation}
\end{theorem}

\begin{proof}
Use row deficiency $j$ instead of rank, so that the kernel has dimension
$k=c+j$.  Write
\[
 u_j=u(c+j),\qquad r_j=\frac{R_{b,c+j}}C.
\]
Then $P(j)=u_jr_j$, $V(j)=u_j$, and
\begin{equation}
 \chi^2(P\Vert V)=\sum_{j=0}^{n}u_j(r_j-1)^2.
 \label{symp:eq:chi-sum}
\end{equation}
We show that the term $j=1$ supplies the leading contribution.

First, the explicit uniform rank law gives
\begin{equation}
 u_1=\frac{Q^c}{q-1}(1-Q^n)
         \prod_{i=c+2}^{n+c}(1-Q^i)
     =\frac{Q^c}{q-1}\bigl(1+O_q(Q^c)\bigr).
 \label{symp:eq:u-one}
\end{equation}
From \eqref{symp:eq:fullrank-ratio}, $r_0=1+O_q(Q^n)$.  The successive
ratio formula gives, including when $h=0$,
\begin{equation}
 r_1=r_0\frac{1-Q^h}{1-Q^n},\qquad
 1-r_1=Q^h\bigl(1+O_q(Q^c)\bigr).
 \label{symp:eq:r-one}
\end{equation}
For clarity, the second estimate follows by writing
$r_0/(1-Q^n)=1+O_q(Q^n)$: subtracting its product with $1-Q^h$ from one
leaves $Q^h+O_q(Q^n)$, and $Q^n/Q^h=Q^c$.  At $h=0$ the identity
$r_1=0$ makes the conclusion exact.  Equations
\eqref{symp:eq:u-one} and \eqref{symp:eq:r-one} yield
\begin{equation}
 u_1(1-r_1)^2
 =\frac{Q^{c+2h}}{q-1}\bigl(1+O_q(Q^c)\bigr).
 \label{symp:eq:chi-leading}
\end{equation}
The $j=0$ contribution is $O_q(Q^{2n})$, a relative
$O_q(Q^c)$ error.

For the other terms we need an estimate that remains valid near the support
boundary.  For every $j\ge1$,
\begin{equation}
 0\le1-r_j\le\frac{Q^{h-2(j-1)}}{1-Q^2}.
 \label{symp:eq:likelihood-defect}
\end{equation}
If $c+j\le m$, then $r_j\ge R_{b,c+j}$, and removing the denominators in
the product for $R$ gives
\[
 R_{b,c+j}\ge\prod_{i=0}^{c+j-1}(1-Q^{b-2i}).
\]
Its defect from one is at most the sum of the displayed powers of $Q$,
which is bounded by the right side of
\eqref{symp:eq:likelihood-defect}.  If $c+j>m$, then $r_j=0$ and
$h-2(j-1)\le0$, so the same bound still holds.  The uniform rank formula
also gives
\begin{equation}
 u_j=Q^{j(c+j)}
       \frac{a_n}{a_{n-j}}\frac{a_b}{a_{c+j}}\frac1{a_j}
     \le a_\infty^{-1}Q^{j(c+j)}.
 \label{symp:eq:rank-tail-bound}
\end{equation}
Consequently, for $j\ge2$, the summand in
\eqref{symp:eq:chi-sum}, divided by $Q^{c+2h}$, is at most
\[
 \frac{Q^{(j-1)c+(j-2)^2}}{a_\infty(1-Q^2)^2}.
\]
Their sum is bounded by
\begin{equation}
 \frac{1}{a_\infty(1-Q^2)^2}
       \frac{Q^c}{1-Q^{c+1}},
 \label{symp:eq:chi-remainder}
\end{equation}
since successive exponents increase by at least $c+1$.  Together with
\eqref{symp:eq:chi-leading}, and $c+2h=2n-c$, this proves
\eqref{symp:eq:chi-square} with uniform error.

For the impossibility assertion, independence gives the exact identity
\[
 1+\chi^2(P^{\otimes t}\Vert V^{\otimes t})
  =\bigl(1+\chi^2(P\Vert V)\bigr)^t.
\]
By Cauchy--Schwarz,
\[
 \operatorname{TV}(P^{\otimes t},V^{\otimes t})
 \le\frac12\sqrt{\bigl(1+\chi^2(P\Vert V)\bigr)^t-1},
\]
which tends to zero when $t=o(q^{2n-c})$.  Optimal equal-prior error is
$(1-\operatorname{TV})/2$, proving the first assertion.

For achievability, let
\[
 p_U=\Pr(\operatorname{rank}U<n),\qquad
 p_A=\Pr(\operatorname{rank}A<n).
\]
The exact sign pattern gives $p_U-p_A=\Delta$.  The uniform full-row-rank
product and a union bound show
\[
 0\le p_A\le p_U
   =1-\prod_{j=c+1}^{n+c}(1-Q^j)
   \le\frac{Q^c}{q-1}.
\]
From $t$ observations compute the fraction failing to have full row rank,
and compare it with $(p_A+p_U)/2$.  Under either hypothesis its variance is
at most $p_U/t$.  Chebyshev's inequality bounds each error probability by
\begin{equation}
 \frac{4p_U}{t\Delta^2}
   =O_q\left(\frac{q^{2n-c}}t\right).
 \label{symp:eq:chebyshev-test}
\end{equation}
This tends to zero in the claimed range.  A sufficiently small constant
multiple of the scale in the chi-square bound, and a sufficiently large
constant multiple in \eqref{symp:eq:chebyshev-test}, give the matching
fixed-error lower and upper bounds.

Finally, the complete-matrix test selects the symplectic model precisely
when the observed kernel is $H$-isotropic.  It never errs under that model,
and errs with probability $\Omega$ under the uniform model.  It is optimal
by \eqref{symp:eq:density}, so its equal-prior error is $\Omega/2$.
Theorem~\ref{symp:thm:separation} gives
\eqref{symp:eq:one-matrix-risk}.
\end{proof}

\subsection{Interpretation and relation to the source}

For example, take any admissible sequence with $c/n\to\alpha\in(0,1)$.
The one-observation rank discrepancy is asymptotic to
$q^{-n}/(q-1)$, and the rank-only sample scale is
$q^{(2-\alpha)n+o(n)}$.  Yet the optimal error from a single complete matrix
is
\[
 \frac12q^{-\alpha^2n^2/2+o(n^2)}.
\]
The information resides in pairings between vectors in the kernel, rather
than in its dimension.  Kernel dimension is almost identically distributed
in the two models; the restriction of the known alternating form to that
kernel separates them.

The supplied article establishes the square case and the general exact
compression law.  The present family keeps the target Lagrangian, takes a
nondegenerate source, and allows a growing excess of columns.  This yields
an explicit solution of one growing-dimension regime, including its
saturation boundary and its rank-only statistical complexity.  It does not
determine the limiting law of arbitrary pairs of growing, possibly
degenerate, restricted forms.  The literature comparison likewise supports
a new extension relative to the supplied corpus; it is not a certification
of worldwide priority.
